\subsection{Limitations}

The results reported in this study are subject to several limitations, most of which stem from the scarcity of country-level plastic waste data and issues regarding their sources.

First, several series we model are extremely short. \texttt{YearTRC (Japan)} and \texttt{YearTRC (United States)} each contain nine annual observations. \texttt{YearPWG (Japan)} contains 26 observations, but its chronological hold-out set contains only two observations. These small samples limit the complexity of models that can be fit reliably and constrain the strength of conclusions about long-run trends. With so few test points, an unusual observation or outlier can dominate the reported metrics, and the resulting model rankings are correspondingly unstable. Therefore, the per-experiment rankings throughout Section~4 should be read as indicative rather than definitive.

Second, some evaluation metrics are unreliable on datasets of this size. In particular, the denominator of \(R^2\) becomes very small when only two observations are held out, so that modest absolute errors can produce large negative values that are difficult to interpret. For this reason, we omit \(R^2\) from the smallest tasks and rely primarily on error-based metrics such as WAPE.

Third, this study does not account for the organizations and reports from which individual data points are derived. Because different reports may adopt different methodologies, techniques, and definitions, and because environmental data are inherently noisy, the pooled heterogeneous data exhibit serious inconsistencies. The correlation and feature-importance analyses are descriptive rather than causal. We therefore treat this study as a preliminary investigation into data-driven modeling of plastic waste, and we leave the development of models that can accommodate these inconsistencies to future work.

Finally, there is a boundary on external reproducibility. All experiments in this study are fully reproducible from the frozen, derived country--year datasets that we release. However, the upstream data pipeline is not, since the PRISM platform and its underlying source records are subject to access restrictions and are not publicly available, so the extraction and validation of records from the original heterogeneous documents cannot be independently re-run by external researchers.
